Las múltiples aplicaciones de los sistemas VIO y VI-SLAM han impulsado el desarrollo de soluciones en diversas ramas de la ciencia, orientadas a abordar el problema en entornos muy variados. A continuación se revisan los trabajos más relevantes en este campo.

Los primeros enfoques para odometría visual-inercial se apoyaron en la teoría de filtros de Kalman para estimar la pose de un dispositivo. MSCKF~\cite{mourikis2007msckf} formuló el problema como un filtro de Kalman extendido (EKF), cuya contribución principal fue excluir las posiciones 3D de los puntos del vector de estado, logrando una complejidad computacional lineal en el número de \textit{features} observados y alcanzando precisión competitiva en escenarios reales de gran escala. ROVIO~\cite{bloesch2015rovio} adoptó un enfoque directo mediante un filtro de Kalman extendido iterativo (IEKF), empleando el error fotométrico sobre parches de imagen como medición. La detección de keypoints se realiza con FAST~\cite{rosten2006fast} y el tracking opera sobre parches multi-nivel.

A diferencia de los métodos basados en filtros, la optimización de mínimos cuadrados no lineales demostró ofrecer mayor precisión y eficiencia computacional. Tras consolidarse en sistemas de SLAM visual como PTAM~\cite{klein2007ptam} y ORB-SLAM~\cite{mur2015orbslam}, se estableció como el paradigma predominante en VI-SLAM. Estos sistemas seleccionan un conjunto reducido de frames (keyframes) mediante una ventana deslizante y realizan alguna variante de bundle adjustment sobre ellos. El desarrollo de este enfoque se ha visto impulsado por librerías de optimización de código abierto como g2o~\cite{kummerle2011g2o}, Ceres Solver~\cite{agarwal2023ceres} y GTSAM~\cite{dellaert2012gtsam}.

OKVIS~\cite{leutenegger2015okvis} fue el primer sistema en formular la fusión visual-inercial como un problema de mínimos cuadrados no lineales en ventana deslizante, acoplando errores de reproyección y preintegración de la IMU.

Basalt~\cite{usenko2020basalt} es un método de VIO basado en optimización de keyframes, dividido en dos componentes: el \textit{front-end} y el \textit{back-end}. El primero emplea FAST~\cite{rosten2006fast} para la detección de features y \textit{optical flow} por KLT~\cite{lucas1981klt} para el tracking. El segundo optimiza una ventana de keyframes con una heurística de marginalización que preserva la eficiencia computacional: cada vez que un keyframe es eliminado de la ventana, se calcula un término de marginalización que condensa su información, garantizando que su contribución se propague a las optimizaciones posteriores. Adicionalmente, los datos marginalizados por el VIO son procesados de forma \textit{offline} por un módulo separado que construye un mapa global, recupera nuevas correspondencias entre puntos y las refina mediante bundle adjustment global sobre todas las poses y landmarks. Cabe destacar que la optimización local de bundle adjustment emplea una formulación de factorización QR~\cite{demmel2022slidingrootba}, adaptada de RootBA~\cite{demmel2021largescalerootba}, que permite operar en precisión simple con una mejora sustancial de rendimiento. Este conjunto de decisiones de diseño le otorga a Basalt una precisión comparable al estado del arte y una elevada robustez, convirtiéndolo en una alternativa sólida para escenarios de VR/AR.

Kimera~\cite{rosinol2020kimera} es una librería modular que integra componentes especializados para el front-end, el back-end y el loop closure mediante pose graph optimization robusto. Como aporte adicional, incorpora un módulo capaz de generar una reconstrucción de malla con etiquetado semántico en 3D.

ORB-SLAM3~\cite{campos2021orbslam3} unificó múltiples modalidades de entrada con soporte para IMU, asociación de datos a corto, mediano y largo plazo, y un sistema de mapas Atlas para gestionar mapas desconectados. Su arquitectura se organiza en tres threads principales: tracking, \textit{local mapping} y \textit{loop closing}. Este último es responsable de detectar regiones comunes entre el mapa activo y los restantes, cerrar ciclos o fusionar mapas según el caso, y lanzar un proceso de bundle adjustment global cuando la cantidad de keyframes lo permite.

El procedimiento de loop closure comienza con la búsqueda de keyframes similares en una base de datos DBoW2~\cite{galvezlopez2012dbow2}. Para cada candidato identificado, se establecen correspondencias 2D-2D y 3D-3D dentro de una ventana local de keyframes covisibles, y se estima la transformación relativa entre el keyframe más reciente y el candidato mediante RANSAC con el algoritmo de Horn~\cite{horn1987algorithm}, seguido de una optimización no lineal. Una vez validado el ciclo, la corrección se realiza en dos fases. En la primera, se forman dos ventanas de keyframes: una centrada en el keyframe más reciente y sus covisibles, y otra en el candidato detectado; los keyframes y landmarks de la primera se actualizan con la transformación estimada y se ejecuta un bundle adjustment local que integra la información de ambas. En la segunda fase, se aplica pose graph optimization sobre el \textit{Essential Graph}, un grafo construido sobre un árbol de expansión mínimo, aristas de loop closures previos y aristas de alto peso del grafo de covisibilidad, propagando la corrección a todos los keyframes del mapa y garantizando la consistencia global.

OKVIS2~\cite{leutenegger2022okvis2} extiende OKVIS con un mapa persistente y mecanismos eficientes de loop closure y marginalización, manteniéndose entre los sistemas de mayor precisión del estado del arte. Su arquitectura comprende un front-end y un \textit{realtime estimator} que procesan de forma sincrónica las imágenes y las mediciones de la IMU. El front-end incorpora la capacidad de segmentar cada frame mediante una red neuronal convolucional (CNN) para identificar áreas dinámicas. Por su parte, el realtime estimator optimiza el grafo de factores, marginaliza observaciones antiguas, fija estados pasados y mantiene un grafo de poses que alimenta la etapa de loop closure.

El proceso de loop closure se activa cuando la base de datos DBoW2 devuelve un candidato válido y la verificación geométrica, mediante RANSAC y el algoritmo gP3P~\cite{kneip2013gpnp} con correspondencias 3D-2D, resulta exitosa. En ese caso, la ventana de keyframes activa se alinea con la pose del keyframe candidato y se fusionan los landmarks redundantes. Para corregir la inconsistencia acumulada, se aplica primero \textit{rotation averaging} y una redistribución uniforme del error de posición entre los keyframes involucrados, y luego se ejecuta una optimización global sobre una copia del mapa completo, fijando los estados ajenos al ciclo e incluyendo los términos de error de la IMU, con lo que se asegura la coherencia global del mapa.

Snake-SLAM~\cite{ruckert2021snakeslam} es un sistema VI-SLAM que también logra métricas de precisión del estado del arte. Su back-end propone una optimización desacoplada que resuelve el estado de la IMU y el bundle adjustment visual como subproblemas independientes, reduciendo la complejidad computacional sin sacrificar precisión. Esto le permite operar con ventanas locales de keyframes más amplias que otros sistemas de SLAM.

Basalt-XR~\cite{zanotti2024basaltxr} extiende Basalt al contexto de VR/AR, añadiendo soporte para más de dos cámaras no paralelas y tracking de features asistido por IMU. Incorpora además un mapa persistente de keyframes y landmarks que habilita una consulta de covisibilidad: el back-end puede reintroducir keyframes del mapa en la ventana local, y el front-end puede reidentificar puntos mediante un proceso de reproyección denominado \textit{recall}.

Una línea reciente de trabajos ha explorado el uso de redes neuronales en el pipeline de SLAM, alcanzando precisión del estado del arte. DROID-SLAM~\cite{teed2021droid} y ViSTA-SLAM~\cite{zhang2026vistaslam} estiman la pose mediante redes neuronales y generan reconstrucciones densas de puntos 3D, mientras que Deep Patch Visual SLAM~\cite{lipson2024dpvslam} produce una reconstrucción sparse; ViSTA-SLAM, además, prescinde de los parámetros intrínsecos de la cámara. Una vertiente complementaria persigue además representaciones fotorrealistas del entorno: Photo-SLAM~\cite{huang2024photoslam} y HI-SLAM2~\cite{zhang2025hislam2} emplean 3D Gaussian splatting como representación del mapa para lograr mapeo fotorrealista y reconstrucción densa, mientras que VIGS-SLAM~\cite{zhu2026vigsslam} suma información inercial para ganar robustez. Si bien prometedores, todos estos sistemas comparten una elevada demanda computacional y la dependencia de GPUs, lo que limita su aplicabilidad en plataformas con recursos acotados, objetivo central del presente trabajo.

En síntesis, OKVIS2 y ORB-SLAM3 son los sistemas más precisos~\cite{krishnan2025lamaria,leutenegger2022okvis2}, pero no alcanzan robustez u operación en tiempo real en entornos de VR/AR~\cite{demayo2025monadoslamdataset}. Los métodos basados en aprendizaje, si bien prometedores, requieren mayor capacidad de cómputo y entrenamiento supervisado. Basalt, en contraste, equilibra robustez, eficiencia y precisión competitiva, lo que lo convierte en la base más adecuada para el sistema propuesto en este trabajo.
